% Consolidated 04-gravitation from b597304e78a782cd65c7976ce94a1a744dfe8067; original TeX fragments preserved.

% === BEGIN 04-gravitation-200-210.tex blob 2da899b248aae0ce591c8623c2a38917bcb329b7 ===
% Irodov, user-supplied Problems in General Physics, printed pp.33--34
% and continuations on printed pp.34--35, source PDF pp.42--44.
% In this chapter G is the gravitational constant and M_s is the solar mass.
\exproblem{200}{33}
一颗行星在太阳引力作用下作圆周运动，日心系速率为
$v=34.9\,\mathrm{km/s}$。求该行星的公转周期。
\exsolution
记太阳质量 $M_\odot$，引力参数
$\mu=GM_\odot$。圆轨道上的引力提供向心加速度：
\[
\frac{v^2}{r}=\frac{\mu}{r^2}
\Longrightarrow r=\frac{\mu}{v^2}.
\]
公转周期 $T=2\pi r/v$，故
\[
\boxed{T=\frac{2\pi GM_\odot}{v^3}}.
\]
取 $\mu\simeq1.3271\times10^{20}\,\mathrm{m^3/s^2}$、
$v=3.49\times10^4\,\mathrm{m/s}$，
算得 $T\simeq2.27\times10^2\,\mathrm d$（约 $227$ 天）。
\excheck 原书简答给出约 $225$ 天，属于所取太阳常数与
有效数字的近似差别；周期计算必须有 $v^3$ 在分母，
否则量纲不为时间。

\exproblem{201}{33}
木星绕太阳的周期是地球的 $12$ 倍，视二者轨道为圆。
求：（a）木星日心距离是地球的几倍；
（b）木星在日心系中的速率、向心加速度。
\exsolution
开普勒第三定律 $T^2\propto r^3$，
故
\[
\boxed{\frac{r_J}{r_E}=12^{2/3}\simeq5.24}.
\]
由圆轨道速率 $v=2\pi r/T$ 与
$a=v^2/r$，
\[
\frac{v_J}{v_E}=12^{-1/3},\qquad
\frac{a_J}{a_E}=12^{-4/3}.
\]
取地球公转速率 $v_E\simeq29.8\,\mathrm{km/s}$、
向心加速度 $a_E\simeq5.93\times10^{-3}\,\mathrm{m/s^2}$：
\[
\boxed{v_J\simeq13.0\,\mathrm{km/s},\qquad
a_J\simeq2.16\times10^{-4}\,\mathrm{m/s^2}}.
\]
\excheck 木星轨道半径较大、速率和向心加速度均较小；
这里用的是日心近圆轨道近似，不是木星相对地球的速度。

\exproblem{202}{33}
行星绕太阳沿椭圆轨道运动，
其最近、最远日心距离分别为 $r$ 和 $R$。
利用开普勒定律求公转周期。
\exsolution
椭圆的半长轴
$a=(r+R)/2$。
按开普勒第三定律，与同一太阳引力参数下
半径为 $a$ 的圆轨道具有同样的周期：
\[
T^2=\frac{4\pi^2a^3}{GM_\odot}.
\]
因此
\[
\boxed{T=2\pi\sqrt{\frac{[(r+R)/2]^3}{GM_\odot}}
=\pi\sqrt{\frac{(r+R)^3}{2GM_\odot}}}.
\]
\excheck 当 $R=r$，表达式退化为
$T=2\pi\sqrt{r^3/(GM_\odot)}$，
恰为圆轨道公式。

\exproblem{203}{33--34}
一小天体在日心参考系内，从与地球公转轨道半径相等
的距离 $R_E$ 处由静止开始落向太阳。
利用开普勒定律估算下落时间，忽略太阳表面半径。
\exsolution
径直下落可看作离心率趋近 $1$ 的退化椭圆轨道的半程：
远日点为 $R_E$，近日点趋于 $0$，半长轴为 $R_E/2$。
其完整退化椭圆的形式周期与同半长轴的圆轨道相同：
\[
T_*=\frac{T_E}{2^{3/2}},
\]
其中 $T_E$ 是地球一年公转周期。
从远日点落到近日点只经过半个周期，故
\[
\boxed{t_{\rm fall}=\frac{T_E}{4\sqrt2}
\simeq64.6\,\mathrm d\approx65\,\mathrm d}.
\]
\excheck 忽略了太阳有限半径；若以太阳表面为终点，
实际下落时间会略短于该中心点极限。

\exproblem{204}{34}
把太阳系所有线性尺寸按同一比例 $\eta$ 缩小，
所用材料的平均密度仍与原恒星、行星各自的平均密度相同。
问行星模型公转周期相对原系统如何变化？
\exsolution
设所有长度缩小为原来的 $1/\eta$
（若 $\eta>1$）。平均密度不变意味着各质量按
体积缩放，故中心天体质量
$M'_\odot=M_\odot/\eta^3$。
行星轨道半长轴 $a'=a/\eta$，
开普勒第三定律给出
\[
\frac{{T'}^2}{T^2}
=\frac{(a')^3/M'_\odot}{a^3/M_\odot}
=\frac{(a^3/\eta^3)/(M_\odot/\eta^3)}
{a^3/M_\odot}=1.
\]
因此
\[
\boxed{T'=T}.
\]
\excheck 系统公转速率按线性尺寸同比缩小，
但周期在相同平均密度的相似模型中保持不变。

\exproblem{205}{34}
一对双星在万有引力下彼此绕共同质心作圆周运动。
已知两星总质量为 $M$、公转周期为 $T$，
求两星之间的距离 $\ell$。
\exsolution
设两星质量为 $m_1,m_2$，间距为 $\ell$，
共同角速度 $\omega=2\pi/T$。
关于质心的轨道半径分别为
$r_1=m_2\ell/M$、$r_2=m_1\ell/M$。
对第一星列向心力方程：
\[
\frac{Gm_1m_2}{\ell^2}
=m_1\omega^2\frac{m_2}{M}\ell.
\]
约去质量因子得
$\omega^2=GM/\ell^3$，所以
\[
\boxed{\ell=\left(\frac{GMT^2}{4\pi^2}\right)^{1/3}}.
\]
\excheck $M=m_1+m_2$ 是两星\emph{总质量}，
并非任意一颗星的质量；结论与质量比无关。

\exproblem{206}{34}
求引力相互作用势能：（a）两个质量分别为 $m_1,m_2$、
间距 $r$ 的质点；（b）质量为 $m$ 的质点位于
一均匀细杆轴线的延长线上，距杆近端为 $a$，
细杆质量 $M$、长度 $\ell$。
再求（b）的引力大小。
\exsolution
取无穷远相互作用势能为零。
（a）按牛顿万有引力积分：
\[
\boxed{U_a(r)=-\frac{Gm_1m_2}{r}}.
\]
（b）均匀细杆线密度 $\lambda=M/\ell$，
从杆近端起坐标 $x\in[0,\ell]$，
微元质量为 $\mathrm dm=\lambda\,\mathrm dx$。
质点与微元距离为 $a+x$，故
\[
U_b=-Gm\int_0^\ell\frac{\lambda\,\mathrm dx}{a+x}
=\boxed{-\frac{GmM}{\ell}
\ln\left(1+\frac{\ell}{a}\right)}.
\]
朝向细杆的引力大小由势能对间隔的导数得到：
\[
\boxed{F=Gm\int_0^\ell
\frac{\lambda\,\mathrm dx}{(a+x)^2}
=\frac{GmM}{a(a+\ell)}}.
\]
\excheck $\ell\to0$ 而总质量 $M$ 不变时，
$U_b\to-GmM/a$，$F\to GmM/a^2$，
退化为两个质点。

\exproblem{207}{34}
一质量为 $m$ 的行星绕太阳作椭圆轨道，
最远、最近日心距离分别为 $r_1,r_2$。
求相对太阳中心的轨道角动量大小。
\exsolution
记太阳引力参数 $\mu=GM_\odot$，
近、远日点切向速率分别为 $v_2,v_1$。
角动量守恒：
$L=mr_1v_1=mr_2v_2$。
机械能守恒：
\[
\frac{L^2}{2mr_1^2}-\frac{GmM_\odot}{r_1}
=\frac{L^2}{2mr_2^2}-\frac{GmM_\odot}{r_2}.
\]
移项并约分可得
$L^2/(2m)=GmM_\odot r_1r_2/(r_1+r_2)$，
因此
\[
\boxed{L=m\sqrt{\frac{2GM_\odot r_1r_2}{r_1+r_2}}}.
\]
\excheck 若 $r_1=r_2=r$，公式化为
$L=m\sqrt{GM_\odot r}$，正是圆轨道角动量。

\exproblem{208}{34}
利用守恒定律证明：质量为 $m$ 的行星绕太阳作椭圆轨道时，
其总机械能仅由椭圆半长轴 $a$ 决定；求 $E(a)$。
\exsolution
仍记 $\mu=GM_\odot$，最近、最远距离
分别为 $r_p,r_a$，所以 $a=(r_p+r_a)/2$。
在两个拱点径向速度均为零，角动量 $L$
及总机械能相同：
\[
E=\frac{L^2}{2mr_p^2}-\frac{\mu m}{r_p}
=\frac{L^2}{2mr_a^2}-\frac{\mu m}{r_a}.
\]
先消去 $E$ 可得
$L^2/(2m)=\mu m r_pr_a/(r_p+r_a)$；
代入任一拱点：
\[
E=\frac{\mu m r_a}{r_p(r_p+r_a)}
-\frac{\mu m}{r_p}
=-\frac{\mu m}{r_p+r_a}.
\]
故
\[
\boxed{E=-\frac{GM_\odot m}{2a}}.
\]
\excheck $E$ 为负表示引力束缚轨道，
且只依赖半长轴、不单独依赖偏心率。

\exproblem{209}{34}
行星在椭圆轨道上的某一时刻距离太阳为 $r_0$、
速率为 $v_0$，速度与日心径矢夹角为 $\alpha$。
求这一轨道的最近和最远日心距离。
\exsolution
记 $\mu=GM_\odot$，单位质量的守恒能量
$\varepsilon=v_0^2/2-\mu/r_0<0$，
比角动量 $h=r_0v_0\sin\alpha$。
定义无量纲参数
$q=r_0v_0^2/\mu$，椭圆条件 $0<q<2$。
半长轴由 $\varepsilon=-\mu/(2a)$ 得
$a=r_0/(2-q)$。
偏心率由能量—角动量关系
\[
e^2=1+\frac{2\varepsilon h^2}{\mu^2}
=1-q(2-q)\sin^2\alpha.
\]
所以两个拱点距离为
\[
\boxed{r_{\min,\max}=
\frac{r_0}{2-q}
\left[1\mp\sqrt{1-q(2-q)\sin^2\alpha}\right]}.
\]
\excheck 负号对应近日点，正号对应远日点；
若输入参数不满足束缚轨道条件，不可称远日点有限。

\exproblem{210}{34}
一宇宙天体在距太阳很远处以速率 $v_0$
向太阳接近，其入射渐近直线距太阳中心的冲量参数
（瞄准距离）为 $\ell$（图1.51）。求最近日心距离。
\exsolution
记 $\mu=GM_\odot$，质点质量在比能量与比角动量中相消。
在无穷远处有
$\varepsilon=v_0^2/2$、$h=\ell v_0$。
最近点处径向速度为零，令该距离为 $r_{\min}$。
由能量守恒
\[
\frac{v_0^2}{2}
=\frac{h^2}{2r_{\min}^2}-\frac{\mu}{r_{\min}},
\]
于是
$v_0^2r_{\min}^2+2\mu r_{\min}-\ell^2v_0^2=0$。
取正根：
\[
\boxed{r_{\min}=\frac{\mu}{v_0^2}
\left[\sqrt{1+\left(\frac{\ell v_0^2}{\mu}\right)^2}-1\right]}.
\]
\excheck 当 $\ell=0$，理想质点可直冲中心；
当 $\mu\to0$ 时应以有理化形式取极限，
可得 $r_{\min}\to\ell$，符合自由直线运动。
% === END 04-gravitation-200-210.tex ===

% === BEGIN 04-gravitation-211-220.tex blob ca689eeb2d48306d1df29d8df9a7ad23e773f605 ===
% Irodov supplied PDF pp.43--44, printed pp.34--35.
% Gravitational shell theorem, self-energy, pressure, satellite phase recurrence.
\exproblem{211}{34}
质量为 $m$ 的质点位于均匀球体外，
球体总质量 $M$、半径 $R$，
质点距球心为 $r>R$。
求：（a）引力相互作用势能；（b）球体施加于质点的引力。
\exsolution
先将球体分解为许多同心薄球壳。由球壳定理，
壳外的引力势与把该壳全部质量集中在球心时相同。
将所有壳的质量积分，球外的总引力势
$\Phi(r)=-GM/r$（取无穷远为零）。
因此
\[
\boxed{U(r)=m\Phi(r)=-\frac{GMm}{r}},
\]
而引力矢量由 $-\nabla U$ 给出：
\[
\boxed{\bm F(r)=-\frac{GMm}{r^3}\bm r},
\qquad |\bm F|=\frac{GMm}{r^2}.
\]
\excheck 球体必须具有球对称密度分布；
“球外等效质点”并不适用于任意非球形物体。

\exproblem{212}{34}
证明均匀薄球壳内部任意一点 $A$ 处的引力矢量为零。
\exsolution
设球壳半径 $R$、总质量 $M$，
内部考察点距壳心为 $r<R$。
壳面质量密度为 $\sigma=M/(4\pi R^2)$。
按以球心为极点的极角 $\theta$ 分环，
环质量
$\mathrm dM=2\pi R^2\sigma\sin\theta\,\mathrm d\theta$，
该环上每个质量元到 $A$ 的距离
$d=\sqrt{R^2+r^2-2Rr\cos\theta}$。
引力势直接积分为
\[
\Phi(r)=-G\int\frac{\mathrm dM}{d}
=-\frac{GM}{2}\int_0^\pi
\frac{\sin\theta\,\mathrm d\theta}
{\sqrt{R^2+r^2-2Rr\cos\theta}}.
\]
令 $u=R^2+r^2-2Rr\cos\theta$，积分得
\[
\Phi(r)=-\frac{GM}{2Rr}
\bigl[(R+r)-|R-r|\bigr].
\]
因 $r<R$，有 $|R-r|=R-r$，
所以 $\Phi(r)=-GM/R$ 为常数，
\[
\boxed{\bm g=-\nabla\Phi=\bm0\qquad(r<R)}.
\]
\excheck 球壳内势不等于零，而是常数；
“力为零”来自势的空间梯度为零。

\exproblem{213}{34--35}
质量为 $m$ 的粒子从一个均匀实心半球
平底面的圆心处被移到无穷远。
半球质量为 $M$、半径 $R$。
求\emph{半球引力}在此过程中做的功。
\exsolution
将同样密度的两个半球拼成质量 $2M$ 的均匀完整球，
原位置正好处于该球的中心。
完整球在中心的引力势为
$\Phi_{\rm full}(0)=-3G(2M)/(2R)=-3GM/R$。
上下两个相同半球对该点势的贡献相等，
因此半球在平底面圆心处产生的势为
$\Phi_{\rm hemi}=-3GM/(2R)$。
初始势能 $U_i=m\Phi_{\rm hemi}$，无穷远 $U_f=0$。
保守引力做功为势能减少：
\[
\boxed{W_g=U_i-U_f=-\frac{3GMm}{2R}}.
\]
\excheck 将粒子搬到无穷远时外力须做正功，
但题目所问的是\emph{引力本身做功}，故答案为负。

\exproblem{214}{35}
质量为 $M$、半径为 $R$ 的均匀实心球，
求球心距离 $r$ 处引力场强矢量 $\bm g(r)$
及以无穷远为零点的引力势 $\Phi(r)$，
分球内、球外给出，并说明随 $r$ 的图形变化。
\exsolution
由球壳定理，半径 $r<R$ 内部包含的质量
$M(r)=M r^3/R^3$，
故球内场强
$\bm g=-GM\bm r/R^3$。
球外等效为中心质点，得到
\[
\boxed{\bm g(\bm r)=
\begin{cases}
-\dfrac{GM}{R^3}\bm r,&0\le r\le R,\\[4pt]
-\dfrac{GM}{r^3}\bm r,&r\ge R .
\end{cases}}
\]
球外势为 $\Phi=-GM/r$。
球内利用
$\Phi(r)=\Phi(R)-\int_R^r g_s(s)\,\mathrm ds$，
其中带符号径向场强 $g_s(s)=-GMs/R^3$，可得
\[
\boxed{\Phi(r)=
\begin{cases}
-\dfrac{GM}{2R}\left(3-\dfrac{r^2}{R^2}\right),
 &0\le r\le R,\\[4pt]
-\dfrac{GM}{r},&r\ge R .
\end{cases}}
\]
\excheck $r=R$ 时场强、势均连续；
场强模从球心零线性增长到表面 $GM/R^2$，随后
按 $r^{-2}$ 减小；势从中心 $-3GM/(2R)$ 单调升至零。

\exproblem{215}{35}
均匀大球密度为 $\rho$，
其中挖去一个球形空腔，空腔球心
相对原大球球心的矢量为 $\bm \ell$。
求空腔内部任一点的引力场强。
\exsolution
使用叠加法：有空腔的球等于完整的密度为 $\rho$
的大球加上一个空腔位置处密度为 $-\rho$ 的虚拟小球。
选原大球球心为坐标原点，空腔内一点位置为 $\bm r$。
完整球在该点的场强为
$-(4\pi G\rho/3)\bm r$。
虚拟小球在该点的场强为
$+(4\pi G\rho/3)(\bm r-\bm\ell)$。
相加得
\[
\boxed{\bm g_{\rm cavity}=-\frac{4\pi G\rho}{3}\bm\ell}.
\]
\excheck 场强与空腔内考察点的具体位置无关，是均匀场；
若空腔与大球同心（$\bm\ell=\bm0$），场强为零。

\exproblem{216}{35}
质量 $M$、半径 $R$ 的均匀自引力球体，
求自引力压缩造成的内部压强 $p(r)$，
并把地球理想化成均匀球估算地心压强。
\exsolution
设密度 $\rho=3M/(4\pi R^3)$。
球内引力场强模为
$g(r)=GM r/R^3=(4\pi/3)G\rho r$。
静力学平衡微分方程
$\mathrm dp/\mathrm dr=-\rho g(r)$，
表面边界条件 $p(R)=0$。积分得
\[
p(r)=\int_r^R\frac{4\pi}{3}G\rho^2s\,\mathrm ds
=\frac{2\pi}{3}G\rho^2(R^2-r^2),
\]
亦即
\[
\boxed{p(r)=\frac{3GM^2}{8\pi R^4}
\left(1-\frac{r^2}{R^2}\right)}.
\]
地心 $p(0)=3GM^2/(8\pi R^4)$。
用地球 $M\simeq5.97\times10^{24}\,\mathrm{kg}$、
$R\simeq6.37\times10^6\,\mathrm m$，
得约 $1.7\times10^{11}\,\mathrm{Pa}$
（约 $1.7\times10^6$ 个标准大气压）。
\excheck 这只是\emph{均匀密度}地球的估值；
实际地球密度随深度变化，不应将此直接称为真实地心压强。

\exproblem{217}{35}
求构成以下物体的全部质量之间的\emph{自身引力势能}：
（a）质量 $m$、半径 $R$ 的均匀薄球壳；
（b）质量 $m$、半径 $R$ 的均匀实心球。
\exsolution
（a）均匀球壳面上的引力势为
$\Phi_{\rm shell}=-Gm/R$。
将球壳微元逐一相互作用的势能总和写为
$U=\tfrac12\int\Phi\,\mathrm dm$，
系数 $1/2$ 避免每一对质量元重复计算：
\[
\boxed{U_{\rm shell}=-\frac{Gm^2}{2R}}.
\]
（b）把均匀实心球由半径 $r$ 的内球逐层组装。
密度 $\rho=3m/(4\pi R^3)$，
已装好的内球质量 $M(r)=(4\pi/3)\rho r^3$，
新外壳质量 $\mathrm dm=4\pi\rho r^2\,\mathrm dr$。
装配微元势能
\[
\mathrm dU=-\frac{GM(r)\,\mathrm dm}{r}
=-\frac{16\pi^2G\rho^2}{3}r^4\,\mathrm dr.
\]
积分至 $R$ 并代入密度得
\[
\boxed{U_{\rm solid}=-\frac{3Gm^2}{5R}}.
\]
\excheck 两项都为负；均匀实心球相对相同质量和外半径
的薄壳更紧密，因此自身引力势能更负。

\exproblem{218}{35}
两颗地球卫星在同一平面内分别沿半径
$r_1=7000\,\mathrm{km}$ 和
$r_2=6930\,\mathrm{km}$ 的圆轨道运动。
求它们周期性达到\emph{最小相互距离}的相邻时间间隔，
讨论同向和逆向运行。
\exsolution
取地球引力参数 $\mu=GM_E$。
两颗卫星的角速度大小分别为
$\omega_1=\sqrt{\mu/r_1^3}$、
$\omega_2=\sqrt{\mu/r_2^3}$，且 $\omega_2>\omega_1$。
当两卫星从地心看角位置相同时，二者相距
$|r_1-r_2|$，即最小距离。
若同向运行，角位置之差的增长速率为
$\omega_2-\omega_1$，故
\[
\boxed{\Delta t_{\rm same}=
\frac{2\pi}{\omega_2-\omega_1}
\simeq4.44\,\mathrm d}.
\]
若两卫星逆向运行，相对角速度为
$\omega_1+\omega_2$，故
\[
\boxed{\Delta t_{\rm opposite}=
\frac{2\pi}{\omega_1+\omega_2}
\simeq0.803\,\mathrm h}.
\]
\excheck 原题未指定两卫星旋转方向，因此必须报告两种情况。
数值取 $\mu\simeq3.986\times10^{14}\,\mathrm{m^3/s^2}$；
均为最小距离事件的\emph{重复间隔}，并非各自公转周期。

\exproblem{219}{35}
比较三种加速度的大小之比：
$g_1$ 为地球表面地球重力所致加速度，
$g_2$ 为地球赤道自转的离心惯性加速度，
$g_3$ 为太阳\emph{直接}对地面物体的万有引力加速度。
\exsolution
取地球表面重力
$g_1\simeq9.81\,\mathrm{m/s^2}$；
赤道处自转离心加速度
$g_2=\Omega_E^2R_E$，
其中 $\Omega_E\simeq7.292\times10^{-5}\,\mathrm{s^{-1}}$、
$R_E\simeq6.37\times10^6\,\mathrm m$，
得到 $g_2\simeq0.0339\,\mathrm{m/s^2}$。
太阳直接引力加速度为
\[
g_3=\frac{GM_\odot}{D^2}
\simeq\frac{1.327\times10^{20}}
{(1.496\times10^{11})^2}
\simeq5.93\times10^{-3}\,\mathrm{m/s^2},
\]
其中 $D$ 为日地距离。因此
\[
\boxed{g_1:g_2:g_3
\simeq1:0.00345:0.000605}.
\]
\excheck $g_3$ 是\emph{太阳的直接引力加速度}，
并非潮汐差分加速度；若误用
$GM_\odot R_E/D^3$ 会低估约 $D/R_E$ 倍。

\exproblem{220}{35}
沿地球极点上空（忽略地球自转影响），
求自由落体加速度比地面值分别减少
$1\%$ 和减少一半时的高度。
\exsolution
设地球半径为 $R_E$，极点地面附近
$g_0=GM_E/R_E^2$。离地高度为 $h$ 时
\[
\frac{g(h)}{g_0}
=\left(\frac{R_E}{R_E+h}\right)^2.
\]
若加速度降为原来的比例 $q$，
解出高度的通式
\[
\boxed{h(q)=R_E(q^{-1/2}-1)}.
\]
对 $q=0.99$ 与 $q=0.50$，分别有
\[
\boxed{h_{1\%}=R_E(1/\sqrt{0.99}-1)
\simeq32\,\mathrm{km}},
\]
\[
\boxed{h_{1/2}=R_E(\sqrt2-1)
\simeq2.64\times10^3\,\mathrm{km}}.
\]
\excheck 题目所给“减少一半”指剩余加速度为 $g_0/2$，
而不是减少 $50\%$ 的高度；地球应近似为球对称。
% === END 04-gravitation-211-220.tex ===

% === BEGIN 04-gravitation-221-233.tex blob ce8518af7cc365d8517fd438f3f4e5b4b15bf892 ===
% Irodov user-supplied Arihant PDF, printed pp.35--36, PDF pp.44--45.
% Final subsection 1.4: terrestrial and lunar satellites and cosmic velocities.
\exproblem{221}{35}
在地球极点使一物体获得竖直向上的初速 $v_0$，
忽略空气阻力。若已知地球半径 $R$ 与地表重力加速度
$g$，求物体可上升的最大高度 $h$。
\exsolution
取地球为球对称天体，记引力参数
$\mu=GM_E=gR^2$。
地表初态的单位质量机械能为
$\varepsilon_i=v_0^2/2-\mu/R$。
如果初速小于逃逸速度，最高点速率为零、半径为 $R+h$，
故机械能守恒：
\[
\frac{v_0^2}{2}-\frac{\mu}{R}
=-\frac{\mu}{R+h}.
\]
移项可得
\[
\boxed{h=\frac{Rv_0^2}{2gR-v_0^2}}.
\]
\excheck 当 $v_0\to0$，$h\simeq v_0^2/(2g)$；
当 $v_0\to\sqrt{2gR}^{\,-}$，$h\to\infty$。
若 $v_0\ge\sqrt{2gR}$，物体在此理想模型中不再有有限最高点。

\exproblem{222}{35}
一颗人造卫星绕地球作圆周运动。
相对不随地球自转的惯性参考系，它的圆轨道速率为 $v$。
给定地球半径 $R$ 和地表重力加速度 $g$，
求卫星离地高度 $h$。
\exsolution
卫星轨道半径为 $r=R+h$，
由向心运动所需加速度等于引力加速度：
\[
\frac{v^2}{r}=\frac{GM_E}{r^2}
=\frac{gR^2}{r^2}.
\]
所以 $r=gR^2/v^2$，从而
\[
\boxed{h=R\left(\frac{gR}{v^2}-1\right)}.
\]
\excheck 近地表圆轨道对应 $v=\sqrt{gR}$、$h=0$；
高空卫星速率较小。这里的 $v$ 不是相对自转地面的速率。

\exproblem{223}{35}
求在地球赤道上空相对地面静止的
\emph{地球同步静止卫星}轨道半径，
以及它相对于地心惯性系的速率和向心加速度。
\exsolution
卫星必须在赤道平面内与地球同向转动，
其角速度等于地球自转角速度
$\Omega_E=2\pi/T_{\rm sidereal}$，
其中恒星日 $T_{\rm sidereal}\simeq86164\,\mathrm s$。
圆轨道条件
\[
\Omega_E^2r=\frac{GM_E}{r^2},
\]
因此
\[
\boxed{r=\left(\frac{GM_E}{\Omega_E^2}\right)^{1/3}
\simeq4.2164\times10^4\,\mathrm{km}},
\]
\[
\boxed{v=\Omega_Er\simeq3.07\,\mathrm{km/s}},\qquad
\boxed{a=\Omega_E^2r\simeq0.224\,\mathrm{m/s^2}}.
\]
\excheck 给出的 $r$ 是从\emph{地心}量起的半径，
不是离地高度（离地高度约 $3.58\times10^4$ km）；
使用的是恒星日而非恰好24小时的太阳日。

\exproblem{224}{35--36}
一颗卫星在赤道平面内沿半径
$r=2.00\times10^4\,\mathrm{km}$ 的圆轨道自西向东运动。
地面赤道某点每隔 $\tau=11.6\,\mathrm h$
看到卫星再次从头顶经过。
已知引力常数 $G$，求地球质量 $M_E$。
\exsolution
卫星自西向东的角速度 $\omega_s$ 快于地球角速度
$\Omega_E=2\pi/T_E$，因此重复过顶的相对角速度
$\omega_s-\Omega_E=2\pi/\tau$。
故
\[
\omega_s=2\pi\left(\frac1{T_E}+\frac1\tau\right).
\]
根据圆轨道引力与向心条件
$GM_E/r^2=r\omega_s^2$，
\[
\boxed{M_E=\frac{4\pi^2r^3}{G}
\left(\frac1{T_E}+\frac1\tau\right)^2}.
\]
代入 $r=2.00\times10^7\,\mathrm m$、
$\tau=11.6\times3600\,\mathrm s$、
$T_E\simeq86164\,\mathrm s$，
得到 $\boxed{M_E\simeq6.0\times10^{24}\,\mathrm{kg}}$。
\excheck 过顶间隔不是卫星本身的公转周期，
必须减去地球自转角速度后再处理。

\exproblem{225}{36}
一颗卫星沿赤道圆轨道从东向西绕地球运动，
轨道半径 $r=1.00\times10^4\,\mathrm{km}$。
求其在\emph{随地球自转的参考系}中的速率与
加速度大小。
\exsolution
记地心惯性系中的西向圆轨道角速度大小为
$\omega_s=\sqrt{GM_E/r^3}$，
地球自转角速度为 $\Omega_E$（东向）。
在地球旋转参考系中，两者角速度方向相反，
相对西向角速度大小为
$\omega_{\rm rel}=\omega_s+\Omega_E$。
故
\[
\boxed{v_{\rm rel}=r(\omega_s+\Omega_E)
=\sqrt{\frac{GM_E}{r}}+\Omega_Er},
\]
\[
\boxed{a_{\rm rel}=r(\omega_s+\Omega_E)^2}.
\]
取 $GM_E=3.986\times10^{14}\,\mathrm{m^3/s^2}$、
$\Omega_E=7.292\times10^{-5}\,\mathrm s^{-1}$，
得 $v_{\rm rel}\simeq7.04\,\mathrm{km/s}$、
$a_{\rm rel}\simeq4.96\,\mathrm{m/s^2}$。
\excheck 这是旋转坐标系中的\emph{运动学相对加速度}，
并非只有地球引力造成的地心惯性系向心加速度。

\exproblem{226}{36}
两颗卫星在地球赤道附近、紧贴地表的理想圆轨道上运动：
一颗顺地球自转方向，另一颗逆向。
求在\emph{地面旋转参考系}中，
逆向卫星动能是顺向卫星动能的多少倍。
\exsolution
记近地圆轨道地心惯性系速率
$v_c=\sqrt{gR}$、地面赤道线速度
$u_E=\Omega_E R$。
相对地面，顺向与逆向卫星的速率分别为
$v_+=v_c-u_E$、$v_-=v_c+u_E$。
相同卫星质量的动能比即速率平方比：
\[
\boxed{\frac{K_-}{K_+}
=\left(\frac{\sqrt{gR}+\Omega_E R}
{\sqrt{gR}-\Omega_E R}\right)^2
\simeq1.27}.
\]
\excheck 若在地心惯性系，两颗相同半径圆轨道卫星的
动能相等；题目特意要求在地面旋转系比较，
因此出现不同的动能数值。

\exproblem{227}{36}
一颗质量为 $m$ 的月球人造卫星沿接近圆形轨道运动，
初始轨道半径是月球半径 $R_M$ 的 $\eta$ 倍。
卫星受宇宙尘埃的微弱阻力，大小为
$F_d=\alpha v^2$，与速度反向。
若下降非常缓慢、轨道始终可近似视为圆，
求卫星从初始轨道降到月球表面的时间。
\exsolution
记月球引力参数 $\mu=GM_M=g_MR_M^2$。
瞬时圆轨道速率 $v=\sqrt{\mu/r}$，
总机械能为 $E=-\mu m/(2r)$。
阻力耗散功率
\[
-\frac{\mathrm dE}{\mathrm dt}
=F_dv=\alpha v^3=\alpha(\mu/r)^{3/2}.
\]
由 $\mathrm dE/\mathrm dr=\mu m/(2r^2)$，
得到
\[
\frac{\mu m}{2r^2}\frac{\mathrm dr}{\mathrm dt}
=-\alpha\frac{\mu^{3/2}}{r^{3/2}},
\quad
\frac{\mathrm dr}{\mathrm dt}
=-\frac{2\alpha\sqrt\mu}{m}\sqrt r.
\]
从 $r(0)=\eta R_M$ 积分到 $R_M$：
\[
\boxed{t\simeq
\frac{m(\sqrt\eta-1)}{\alpha}
\sqrt{\frac{R_M}{\mu}}
=\frac{m(\sqrt\eta-1)}{\alpha\sqrt{g_MR_M}}}.
\]
\excheck 这是绝热缓慢圆轨道衰减的近似式；
若阻力大到单圈内轨道已明显非圆，需解完整运动方程。

\exproblem{228}{36}
计算月球表面附近的第一宇宙速度（圆轨道速率）
和第二宇宙速度（逃逸速率），并与地球对应数值比较。
\exsolution
对于质量 $M$、半径 $R$ 的球对称天体，
第一宇宙速度来自圆轨道条件
$v_1^2/R=GM/R^2$，故
$v_1=\sqrt{GM/R}=\sqrt{gR}$。
第二宇宙速度来自逃逸总机械能为零：
$v_2=\sqrt{2GM/R}=\sqrt2\,v_1$。
月球取 $g_M\simeq1.62\,\mathrm{m/s^2}$、
$R_M\simeq1.737\times10^6\,\mathrm m$，得
\[
\boxed{v_{1,M}\simeq1.68\,\mathrm{km/s},\qquad
v_{2,M}\simeq2.37\,\mathrm{km/s}}.
\]
相应地球理想值约为
$7.9\,\mathrm{km/s}$ 和 $11.2\,\mathrm{km/s}$。
\excheck 两组速度均满足 $v_2/v_1=\sqrt2$；
这里忽略了大气与天体自转。

\exproblem{229}{36}
一飞船沿近似掠过月球表面的抛物线轨道接近月球。
在最近点瞬时开启制动发动机，使轨道转为掠月球表面的
圆轨道。求制动造成的飞船速率\emph{有符号变化}。
\exsolution
抛物线轨道在无穷远处比机械能为零，
最近点 $r\simeq R_M$ 的速率为
$v_p=\sqrt{2GM_M/R_M}$。
圆轨道所需速率为
$v_c=\sqrt{GM_M/R_M}$。
因此发动机制动造成的速率变化
\[
\boxed{\Delta v=v_c-v_p
=(1-\sqrt2)\sqrt{\frac{GM_M}{R_M}}
\simeq-0.69\,\mathrm{km/s}}.
\]
\excheck 负号表示减速；若题目问制动速度增量的大小，
应报告 $|\Delta v|\simeq0.69\,\mathrm{km/s}$。
不是飞船速率“增加”。

\exproblem{230}{36}
飞船已在地球表面附近的理想圆轨道上运行。
要摆脱地球引力束缚，沿原速度方向最少还需给予多大的
瞬时速率增量？
\exsolution
近地圆轨道速率为
$v_c=\sqrt{GM_E/R_E}=\sqrt{gR_E}$。
在相同高度达到逃逸抛物线轨道需要速率
$v_e=\sqrt{2GM_E/R_E}=\sqrt2\,v_c$。
因此同向瞬时加速的最小速度增量为
\[
\boxed{\Delta v=(\sqrt2-1)\sqrt{gR_E}
\simeq3.27\,\mathrm{km/s}}.
\]
\excheck 这里未包括发动机损耗、轨道变轨时间、
空气阻力及其他天体引力；理想情况下同向脉冲最节省增量。

\exproblem{231}{36}
已知地球质量是月球质量的 $\eta=81$ 倍，
地月中心距离是地球半径 $R_E$ 的 $n=60$ 倍。
求地月连线上两者引力场强合成为零的点距月球中心多远。
\exsolution
记地月中心距 $D=nR_E$，
零场点距月球为 $x$，则距地球为 $D-x$。
两引力方向相反，大小相等：
\[
\frac{GM_M}{x^2}=\frac{GM_E}{(D-x)^2}
=\frac{G\eta M_M}{(D-x)^2}.
\]
取正平方根：
$(D-x)/x=\sqrt\eta$，解得
\[
\boxed{x=\frac{D}{1+\sqrt\eta}
=\frac{nR_E}{1+\sqrt\eta}
=6R_E\simeq3.82\times10^4\,\mathrm{km}}.
\]
\excheck 这是静止牛顿引力场\emph{零场点}，
不是地月旋转参考系中的完整拉格朗日 $L_1$ 点；
后者还需考虑离心惯性项。

\exproblem{232}{36}
一质量为 $m=2.0\times10^3\,\mathrm{kg}$ 的飞船
从地球表面被运至月球表面。
忽略大气、太阳作用和能量损耗，估计所需的最小外力功。
\exsolution
为估算下限，令飞船始末两端相对各自行星静止，
采用地月近似不动、距离 $D$ 固定的保守引力模型。
初态与末态引力势能近似为
\[
U_i\simeq-\frac{GM_E m}{R_E}
-\frac{GM_Mm}{D-R_E},\quad
U_f\simeq-\frac{GM_E m}{D-R_M}
-\frac{GM_Mm}{R_M}.
\]
理想可逆慢速运输时外力净功等于势能差
$W_{\rm ext,min}=U_f-U_i$。
因为 $D\gg R_E,R_M$，忽略远端引力修正得到
\[
W_{\rm ext,min}\simeq
m\left(\frac{GM_E}{R_E}-\frac{GM_M}{R_M}\right)
=m(g_ER_E-g_MR_M).
\]
用 $g_E\simeq9.81\,\mathrm{m/s^2}$、
$R_E\simeq6.37\times10^6\,\mathrm m$、
$g_M\simeq1.62\,\mathrm{m/s^2}$、
$R_M\simeq1.737\times10^6\,\mathrm m$，
\[
\boxed{W_{\rm ext,min}\sim1.2\times10^{11}\,\mathrm J
\approx1.2\times10^8\,\mathrm{kJ}}.
\]
\excheck 这是一种准静态理想势能预算；
不能将其视为有时间限制、必须安全着陆的真实任务
的发动机燃料能量最小值。原书简答数量级约为
$1.3\times10^8\,\mathrm{kJ}$。

\exproblem{233}{36}
忽略地球自转，估算从地球表面发射物体，
使它最终能逃出\emph{整个太阳系}所需的最小初速，
即第三宇宙速度。
\exsolution
记地表第一宇宙速度
$v_{1,E}=\sqrt{GM_E/R_E}\simeq7.9\,\mathrm{km/s}$，
地球绕太阳的圆轨道速率
$V_E=\sqrt{GM_\odot/D}\simeq29.8\,\mathrm{km/s}$。
日地轨道半径处太阳逃逸速率为 $\sqrt2V_E$。
为使从地球引力井出来的物体在日心系中达到这一速度，
应沿地球公转方向发射，使地球外的双曲线余速
\[
u_\infty=(\sqrt2-1)V_E.
\]
地表初速 $v_3$ 的单位质量地球中心机械能为
$v_3^2/2-GM_E/R_E$，
在地球引力范围外为 $u_\infty^2/2$，
故
\[
\boxed{v_3=
\sqrt{\frac{2GM_E}{R_E}
+(\sqrt2-1)^2V_E^2}
=\sqrt{2v_{1,E}^2+(\sqrt2-1)^2V_E^2}
\simeq16.7\,\mathrm{km/s}}.
\]
\excheck 最优方向为地球公转切向前进方向；
忽略了地球自转、月球引力、其他行星摄动及真实轨道
交接处的细节，属于经典近似估算。
% === END 04-gravitation-221-233.tex ===
